\section{Exact memory under terminal readout contrast}
\label{sec:r74-contrast}

The normal form separates the cost of a permutation module from the positivity
of its preparation and decoder. We now determine when that linear lower bound
is attained, for every finite observable group, and describe the exact memory
transitions as the terminal contrast varies. This is a change of terminal
readout. It differs from the length-dependent background damping of
Section~\ref{sec:r73-finite-certificates}.

Let $b$ be a nonconstant compactly reversible response with finite observable
group $G$, residual cardinality $N$, and real observable module $V_b$. Set
\begin{equation}\label{eq:r74-contrast}
 \pi=\frac1{|G|}\sum_{g\in G}b(g),\qquad
 b_t(g)=\pi+t\{b(g)-\pi\},\qquad 0\le t\le1.
\end{equation}
If $\pi_c=0$, nonnegativity implies $b(g,c)=0$ for every $g$; such an output
can be removed and later restored as a zero decoder column. Thus positivity
of all remaining $\pi_c$ imposes no additional hypothesis. For prepared
unitary commands, \eqref{eq:r74-contrast} is realized by the terminal
measurement $E_c(t)=tE_c+(1-t)\pi_c I$, with the preparation and commands
unchanged.

\begin{theorem}[Attainment of the module bound at small contrast]
\label{thm:r74-low-contrast}
For every finite observable group $G$ there is an explicitly positive $t_0$
such that
\begin{equation}\label{eq:r74-small-contrast}
 \kappa_G(b_t)=\nu_G(V_b),\qquad
 m(b_t)=\min\{N,1+\nu_G(V_b)\}\qquad(0<t\le t_0).
\end{equation}
For every $t>0$, the observable group, module and residual cardinality agree
with those of $b$. At $t=0$ the minimum is one label. The assertion allows
arbitrary finite groups and arbitrary finite terminal outcome sets.
\end{theorem}
\begin{proof}
Every $f\in V_b$ has a suffix-independent total mass
$\ell(f)=\sum_c f(z,c)$, because this holds for its spanning residuals.
Group averaging sends every residual to the same constant response
$\overline b(z)=\pi$. Its image on $V_b$ is therefore $\R\overline b$, and
\[
 V_b=\R\overline b\oplus W,
 \qquad W=\ker\left(\frac1{|G|}\sum_{g\in G}L_g\right).
\]
The module $W$ has no invariant vector, and $\ell(W)=0$.

Choose a $G$-set $X$ of minimum size $k=\nu_G(V_b)$ and an equivariant
injection $j:W\to\R^X$. The permutation inner product is invariant.
Since group averaging commutes with $j$, its image is orthogonal to the
whole invariant subspace, including the uniform row
$p=(1/k,\ldots,1/k)$. In particular its vectors have coordinate sum zero.
Let $F:\R^X\to W$ be orthogonal projection onto $j(W)$ followed by
$j^{-1}$. It is equivariant, $Fj=I_W$, and $F(p)=0$.

Put $v=b-\overline b$, $a=j(v)$, and
\[
 T(x,c)=F(e_x)(\varnothing,c).
\]
The rows of $T$ sum to zero, $pT=0$, and equivariance gives
$aP_gT=b(g)-\pi$. Since $a\ne0$, we may take
\[
 \epsilon=\frac{1}{2k\|a\|_\infty},\qquad
 B=\max_{x,c}\frac{|T(x,c)|}{\pi_c}>0,
 \qquad t_0=\min\{1,\epsilon/B\}.
\]
The row $\lambda=p+\epsilon a$ is a probability row with every entry at
least $1/(2k)$. Since $(t/\epsilon)B\le1$, for $0\le t\le t_0$ the matrix
\[
 D_t(x,c)=\pi_c+(t/\epsilon)T(x,c)
\]
is row-stochastic. Direct multiplication yields
\[
 \lambda P_gD_t
 =\pi+(t/\epsilon)\{pP_gT+\epsilon aP_gT\}
 =b_t(g).
\]
This proves the upper $\kappa_G(b_t)\le k$ with a bounded decoder.

The linear map
$f\mapsto tf+(1-t)\ell(f)\overline b$ is equivariant and invertible for
$t>0$: it is the identity on $\R\overline b$ and multiplication by $t$ on
$W$. It sends each residual of $b$ to the corresponding residual of $b_t$.
It preserves the group kernel and the equality pattern among residuals.
The module lower bound \eqref{eq:r73-module-bounds} now gives equality,
and Theorem~\ref{thm:r73-memory-formula} charges the pure start. At zero
contrast the decoder is constant.
\end{proof}

The construction is an equivariant positivity argument. Symmetric cone
factorizations and coordinate-symmetric lifts have an established theory
\cite{GPT13,KPT10}; ordinary nonnegative rank also has classical barycentric
perturbation results \cite{BocciCarliniRapallo11}. The assertion here is the
exact value $\nu_G(V_b)$ while preserving one common action for every command
word and the prescribed-start charge. It is not a claim that positivity
under mixing or equivariant factorization was previously unknown.

There is a useful distinction at an infinite compact observable group.
The same invertible change of residuals preserves that group for every
$t>0$, so Theorem~\ref{thm:r72-reversible} still excludes every finite
exact stochastic realization. Thus making the terminal probabilities
strictly interior need not make exact memory finite. Terminal readout
contrast, repeated channel damping, and perturbations of an entire
instrument are different operations.

\begin{theorem}[Attained exact contrast thresholds]
\label{thm:r74-contrast-thresholds}
For each integer $k\ge1$ there is $\tau_k\in[0,1]$ such that
\begin{equation}\label{eq:r74-threshold-set}
 \kappa_G(b_t)\le k\quad\Longleftrightarrow\quad 0\le t\le\tau_k.
\end{equation}
These thresholds are nondecreasing in $k$, satisfy $\tau_N=1$, and obey
\[
 k<\nu_G(V_b)\Longrightarrow\tau_k=0,
 \qquad \tau_{\nu_G(V_b)}>0.
\]
The lower memory value is retained at every threshold itself. There are
only finitely many positive-contrast transitions. If all $b(g,c)$ are real
algebraic, all thresholds are effectively computable real algebraic numbers.
When all subgroups of $G$ are normal, they are obtained by finitely many
linear programs; rational tables then have rational thresholds and rational
optimal controllers.
\end{theorem}
\begin{proof}
For a fixed action on at most $k$ points, the variables $\lambda,D,t$ lie
in a compact product of probability simplexes and $[0,1]$. The equations
$\lambda P_gD=b_t(g)$ define a closed set, so its feasible contrasts form
a compact set. If $s>0$ is feasible, the same $\lambda$ and the decoder
$(t/s)D+(1-t/s)\mathbf1\pi$ make every $0\le t\le s$ feasible.
There are only finitely many actions on at most $k$ labelled points;
their feasible sets therefore form a closed initial interval. A single
fixed point realizes $b_0$, including when no positive contrast is feasible.

The residual machine gives $\kappa_G(b_t)\le N$; the remaining inequalities
follow from Theorem~\ref{thm:r74-low-contrast} and its module lower bound.
The exact memory is recovered from
$m(b_t)=\min\{N,1+\kappa_G(b_t)\}$ for $t>0$.
For algebraic data the equations for each action are finite polynomial
constraints over the real algebraic numbers. Classical real quantifier
elimination computes their projected interval and endpoint, as in
Section~\ref{sec:r73-algebraic}. This is a proved finite procedure,
not a reported general quantifier-elimination execution. The final assertion
is the linear reduction and rationality proved in
Theorems~\ref{thm:r74-quotient-lp} and~\ref{thm:r74-quotient-dual} below.
\end{proof}

\subsection{A sharp transition for the retained prepared-qubit experiment}
The bounded-decoder obstruction of
Proposition~\ref{prop:r73-decoder-obstruction} has a complete contrast law.
Let $u=(1,-1,0)$ and
\[
 \beta_t(g)=\frac12+\frac t4uP_gu^{\mathsf T},\qquad g\in S_3,
\]
and let $b_t=(\beta_t,1-\beta_t)$. This is the same prepared-qubit
experiment with terminal effect $t|+\rangle\langle+|+(1-t)I/2$.

\begin{proposition}[An exact qubit readout threshold]
\label{prop:r74-s3-contrast}
For $0<t\le1$, this response has $d=3$, $N=6$, and
\[
 \kappa_{S_3}(b_t)=
 \begin{cases}3,&0<t\le2/3,\\6,&2/3<t\le1,\end{cases}
 \qquad
 m(b_t)=
 \begin{cases}4,&0<t\le2/3,\\6,&2/3<t\le1.\end{cases}
\]
At $t=0$ its memory is one.
\end{proposition}
\begin{proof}
The group, rank and residual count are unchanged for $t>0$ by
Theorem~\ref{thm:r74-low-contrast}. A faithful $S_3$-set of size at most
five has exactly one natural three-point orbit and otherwise only fixed
points or one two-point sign orbit, as in the proof of
Proposition~\ref{prop:r73-decoder-obstruction}. Let $A$ be the initial
mass on its natural orbit. Write its unnormalized initial row and scalar
decoder as
\[
 p=(A/3)\mathbf1^{\mathsf T}+r,
 \qquad D=q\mathbf1+s,
 \qquad r\mathbf1=0,\quad\mathbf1^{\mathsf T}s=0.
\]
The contribution is $Aq+rP_gs$. All remaining orbits contribute a linear
combination of the constant and sign functions. The functions $rP_gs$ and
$uP_gu^{\mathsf T}$ are orthogonal to both: their group sums vanish because
$|S_3|^{-1}\sum_gP_g=\mathbf1\mathbf1^{\mathsf T}/3$, and
$\sum_g\operatorname{sgn}(g)P_g=0$ entrywise. Consequently equality of the
full responses implies
\begin{equation}\label{eq:r74-s3-coefficient}
 rP_gs=(t/4)uP_gu^{\mathsf T}\qquad(g\in S_3).
\end{equation}

Here no unproved matrix-coefficient injectivity is needed. On the row plane
$\mathbf1^\perp$, use the basis $u,(1,1,-2)$. The matrices of the identity,
$(12)$, $(23)$ and the chronological cycle $(123)$ are respectively
\[
 I_2,\quad
 \begin{pmatrix}-1&0\\0&1\end{pmatrix},\quad
 \frac12\begin{pmatrix}1&1\\3&-1\end{pmatrix},\quad
 \frac12\begin{pmatrix}-1&1\\-3&-1\end{pmatrix}.
\]
They span all real two-by-two matrices: the first two span the diagonal
ones, and the latter two have independent off-diagonal parts. Therefore
\eqref{eq:r74-s3-coefficient} holds with every endomorphism of this plane
in place of $P_g$. Equality of the associated rank-one bilinear functionals,
with $t>0$, forces $r=au$, $s=bu^{\mathsf T}$ and $ab=t/4$.
Nonnegativity of $p$ gives $|a|\le A/3$, while the bounded decoder gives
$|b|\le\min(q,1-q)\le1/2$. Hence
\[
 t/4=|ab|\le A/6\le1/6.
\]
No core of size at most five is possible above $t=2/3$.

For $0<t\le2/3$ use the natural three-point action with
$\lambda=(2/3,0,1/3)$ and scalar decoder
$D=\tfrac12\mathbf1+(3t/4)u^{\mathsf T}$. Its entries lie in $[0,1]$
and direct multiplication gives $\beta_t$; the rank bound proves
$\kappa=3$. Above the threshold the six-residual machine supplies the
matching upper. The pure-start formula gives the stated minima and the
constant case is immediate.
\end{proof}
